\section{Locality and strong commutation at isolated types}
\label{sec:locality}

Fix an isolated source type $i$ and its measured type $p(i)$. Let $B_S$ and
$B_T$ denote their boundaries, equipped with compact-invariant probability
measures. These are vertex boundaries; ordinary chamber variables in the
other components will be integrated out without changing notation. Write
$c_i>0$ for the constant in Equation~\eqref{eq:input-19} at this type.
The scale and height analysis of Sections~\ref{sec:scale} and~\ref{sec:heights}
gives a common homogeneous dimension $Q$ on the two sides. Set
\[
 \vartheta=\frac{2}{Q},\qquad
 K_S(x,z)=k_S(x,z)^\vartheta,\qquad
 K_T(v,u)=k_T(v,u)^\vartheta.
\]
In the groups admitted here, $Q>8$, so $0<\vartheta<1/4$. For a source
boundary point $z$, write
\[
 Q_z=(P_z,\log b_z),\qquad B_z=b_z^\vartheta.
\]
Here $Q_z$ is the joint spectral measure of the indicated commuting pair
at one label. Our first objective is strong commutation as $z$ varies.
We will then prove commutation between the two normal copies sharing the
trunk slot.

The inputs from Section~\ref{sec:heights} are the exact ordered
imaginary-power identity~\eqref{eq:input-19}, the isometry~\eqref{eq:input-22},
and the unitary-group intertwining~\eqref{eq:input-23}, including the
own-coordinate density commutation proved there. Each $b_z$ is positive,
has zero kernel, and has finite spectral values. Source transport gives
strong continuity of bounded continuous own-coordinate spectral functions;
when an unbounded operator varies with a label, we use its bounded resolvents.

The kernel construction below has four roles. Its bilinear formula
makes the transported products $K_T(v,P_z)B_z$ depend linearly on a
finite-dimensional source lift, and its inverse-square integrability
supplies a dense common graph domain. On that domain, its fourth-order
polynomial dependence turns a small off-diagonal estimate into exact
locality. Finally, the embedded cone supplies enough commuting evaluation
multipliers to eliminate every positive-order term in the resulting local
finite-jet formulas. Graph cutoffs and individual cores then carry these
exchange identities to strong commutation.

\subsection{A finite-dimensional boundary kernel}

\begin{lemma}[Polynomial boundary lifts]
\label{lem:locality-kernel}
For either boundary $B=B_S,B_T$ there are a finite-dimensional real
vector space $E$, a nondegenerate symmetric bilinear form $[\ ,\ ]$ on
$E$, and a smooth map $n:B\longrightarrow E$ such that
\begin{enumerate}[label=(\roman*)]
\item $n(B)$ spans $E$, and $(r,z)\mapsto r n(z)$ embeds
$(0,\infty)\times B$ smoothly in $E$;
\item $[n(w),n(z)]=K(w,z)$, with $K(w,z)>0$ for $w\ne z$;
\item in a unipotent boundary chart, $n(z)$ is a polynomial of weighted
degree at most four divided by a fixed smooth positive function;
\item for a compact Carnot distance $d$,
\[
 c\,d(w,z)^4\le K(w,z)\le C\,d(w,z)^4;
\]
\item there are positive constants $a_1,a_2,a_{-2}$ such that, for every
$u\in B$,
\[
 \int_B K(v,u)\,dv=a_1,
 \quad \int_B K(v,u)^2\,dv=a_2,
 \quad \int_B K(v,u)^{-2}\,dv=a_{-2}<\infty.
\]
\end{enumerate}
\end{lemma}

\begin{proof}
We spell out the finite-dimensional representation used here.  For
$\operatorname{Sp}(n,1)$, use the irreducible representation of absolute highest
weight $e_1+e_2$ in the root system $C_{n+1}$.  Normalize the split
generator so that its evaluations on the first two standard
coordinates are both one.  The highest restricted weight is then two,
is attained on a single extremal weight, and has multiplicity one.
The lowest restricted weight is minus two.  For $F_{4(-20)}$, use the
representation of highest short root.  In the customary coordinates
with short roots
\[
 \{\pm e_j\}\ \cup\
 \left\{\tfrac12(\pm e_1\pm e_2\pm e_3\pm e_4)\right\},
\]
the split direction has $e_1(H)=2$ and all other $e_j(H)=0$.
Again the extremal evaluations two and minus two are unique.  The
absolute vanishing root system is $B_3$.  These weight assertions also
follow directly by maximizing the displayed linear functional on the
Weyl orbit and its convex hull.

Here are also the determinant weights, which will identify the kernel
exactly. In type $C_{n+1}$ the roots of split value two are
$2e_1,2e_2,e_1+e_2$, and those of split value one are
$e_i\pm e_j$, $i=1,2$, $j=3,\ldots,n+1$. Thus
\[
 (m_\alpha,m_{2\alpha},Q)=(4n-4,3,4n+2),\qquad
 \sum_{\beta(H)>0}\beta=(2n+1)(e_1+e_2).
\]
For $F_{4(-20)}$ the roots of split value two are $e_1$ and
$e_1\pm e_j$, $j=2,3,4$, and those of split value one are
$\tfrac12(e_1\pm e_2\pm e_3\pm e_4)$. Hence
\[
 (m_\alpha,m_{2\alpha},Q)=(8,7,22),\qquad
 \sum_{\beta(H)>0}\beta=11e_1.
\]
In either case the determinant of the positive nilradical has absolute
highest weight $(Q/2)\lambda$, where $\lambda$ is the highest weight
chosen above. The quaternionic groups in the present class have $n\ge2$;
thus $Q\ge10$ in that case and $Q=22$ in the Cayley case.

We justify the required real representation. Let $\kappa$ be split
Chevalley conjugation on the complex Lie algebra and $\sigma$ its given
real-form conjugation. The types $C_{n+1}$ and $F_4$ have no diagram
automorphisms, so the complex automorphism $\sigma\kappa$ is inner,
say $\Ad(h)$. The highest-weight module has a Chevalley real form, with
antilinear conjugation $C$ intertwining $\kappa$. The map $J=\pi(h)C$
therefore commutes with the given real Lie algebra. Schur's lemma gives
$J^2=a\id$. Since $J$ preserves the unique complex top space for the
real split generator, its square on that line is a positive real scalar.
Thus $a>0$, and rescaling $J$ gives $J^2=\id$. Its fixed space is the
required real module $E$. These highest weights lie in the root lattice,
so the representations descend to the adjoint groups used here.

The longest Weyl element is $-1$ in both types. Self-duality therefore
gives a nondegenerate invariant complex bilinear form. Its conjugate by
$J$ is a scalar multiple of itself; rescaling the form makes it $J$-real,
so it restricts to a nondegenerate real pairing on $E$. Split-weight
orthogonality says that the top line pairs only with the bottom line;
nondegeneracy makes that pairing nonzero. The stabilizer of the top line
contains the boundary parabolic and is proper. In rank one the Bruhat
decomposition has only two cells, so this stabilizer is exactly that
parabolic. Its projective orbit is the smoothly embedded boundary.

Choose a maximal-compact-invariant Euclidean norm on $E$. The boundary
is a sphere, $S^{4n-1}$ or $S^{15}$, and is simply connected. Its real
line bundle consequently has a smooth global unit lift $n$. Connectedness
of the group implies that all its action scalars on this lift are positive.
Distinct ordered boundary pairs form one group orbit, so the pairing is
nonzero with constant sign off the diagonal. Change its sign and scale
to make it positive there and one at standard compact opposites. By
irreducibility its transpose is a scalar multiple of itself, with scalar
$1$ or $-1$; positivity on both orders excludes $-1$. The real pairing
is therefore symmetric.

The determinant highest line computed above generates the Cartan
component of $V_\lambda^{\otimes Q/2}$. Its lift is the tensor power
$n(z)^{\otimes Q/2}$. Compact transitivity makes its unit normalization
agree with that of the determinant line, and its invariant pairing is
$[n(w),n(z)]^{Q/2}$. Normalizing at the same opposite pair removes the
single constant. Consequently
\[
 k(w,z)=[n(w),n(z)]^{Q/2},\qquad
 K(w,z)=[n(w),n(z)].
\]
This uses the entire highest-line representation, not only its split
character.

In a chart based at the lowest line, the unnormalized lift is
$\exp(X)n_-$, where $X$ ranges over the graded positive nilradical.
A term raising restricted weight from minus two to at most two has
total graded degree at most four.  Exponentiation is thus a
polynomial of that weighted degree.  Passing to a unit lift divides
by its smooth positive norm.  Write $q(X)=\exp(X)n_-$ for the unnormalized lift, and let
$\delta_t$ multiply degree-one coordinates by $t$ and degree-two
coordinates by $t^2$. Since the bottom weight is $-2$,
\[
 q(\delta_t X)=t^2\pi(a_{\log t})q(X).
\]
Invariance of the pairing therefore makes its two-point kernel
homogeneous of degree four. Its invariance under simultaneous left
translation in the nilradical reduces it to the difference coordinates.  It is
positive away from zero.  Compactness of the homogeneous unit sphere
gives comparison with the fourth power of a homogeneous gauge;
finitely many charts give the asserted compact distance comparison.

The highest-line orbit is an embedded projective orbit.  The unit
lift and its positive radial parameter consequently give the claimed
cone embedding.  Its span is a nonzero invariant subspace, hence all
of $E$.  Compact transitivity makes the three displayed integrals
independent of $u$.  The first two are finite and positive.  For the
last one, a ball of radius $r$ has measure bounded by $Cr^Q$ and the
integrand is bounded by a constant times $d(v,u)^{-8}$.  Dyadic
annuli give
\[
 \int_{d(v,u)<1}K(v,u)^{-2}\,dv
 \le C\sum_{j\ge0}2^{8j}2^{-Qj}<\infty.
\]
Here $Q=4n+2\ge10$ for $\operatorname{Sp}(n,1)$ and $Q=22$ for $F_{4(-20)}$.
\end{proof}

\subsection{The common graph domain}

On the target chamber-variable Hilbert space in
Equation~\eqref{eq:input-22}, define the positive decomposable
operators
\[
 A_z(v)=K_T(v,P_z)B_z.
\]
Their factors commute because they belong to the single $Q_z$
calculus. All integrals in this Section use the fixed compact-invariant
probabilities; no coefficient trace or matrix size is varied.  Write $I$ and $I_S$ for the constant-function injections on
the target and source chamber-variable spaces, respectively.
Equation~\eqref{eq:input-23}, with its frequency rescaled
by $\vartheta$, gives
\begin{equation}
 V A_z^{\mathrm i t}=C_z^{\mathrm i t}V,
 \qquad
 C_z(Z)=c_i^\vartheta K_S(y(Z),z)b_{y(Z)}^{-\vartheta},
 \qquad t\in\mathbb R.
 \label{eq:locality-positive-exchange}
\end{equation}
The source point $y(Z)$ is the vertex of the ordinary source chamber
$Z$ at the opposite type.  In particular, for fixed $Z$ it is
independent of $z$.

\begin{lemma}[Same-variable commutation and norm identities]
\label{lem:locality-graph-identities}
For every $z,w\in B_S$ and $v\in B_T$, the operators $A_z(v)$ and
$A_w(v)$ strongly commute.  The domain
\[
 \mathfrak D
 =\left\{\xi\in\mathcal H:
          \int_{B_S}\|B_z\xi\|^2\,dz<\infty\right\},
 \qquad
 \|\xi\|_{\mathfrak D}^2
 =\|\xi\|^2+\int_{B_S}\|B_z\xi\|^2\,dz
\]
is a Hilbert graph domain.  Its vectors belong to every $\operatorname{dom} B_z$,
and
\begin{equation}
 \sup_z\|B_z\xi\|^2
 \le C\int\|B_x\xi\|^2\,dx,
 \qquad
 \|B_z\xi\|^2+\|B_w\xi\|^2
 \ge c_1 d(z,w)^8\int\|B_x\xi\|^2\,dx.
 \label{eq:input-24}
\end{equation}
On $\mathfrak D$, the maps $z\mapsto A_z(v)$ and $z\mapsto B_z$
are linear in the coordinates of $n_S(z)$.  The map $v\mapsto A_z(v)$
is linear in the coordinates of $n_T(v)$.
\end{lemma}

\begin{proof}
The range of $V$ reduces each unitary group $C_z^{\mathrm i t}$:
Equation~\eqref{eq:locality-positive-exchange} applies to both $t$ and
$-t$.  The $C_z$ strongly commute, since at each ordinary source
chamber they are scalar multiples of the same operator
$b_{y(Z)}^{-\vartheta}$.  The spectral theorem transfers their
resolvents to the target representation.  Hence $A_z(v)$ and $A_w(v)$
strongly commute for almost every $v$.  For fixed $z$, bounded
resolvents of $A_z(v)$ depend strongly continuously on $v$, by bounded
convergence in the $Q_z$ calculus.  Products of their bounded
resolvents have the same continuity.  Their commutation therefore
holds for every $v$.

Positive spectral truncation in
Equation~\eqref{eq:locality-positive-exchange} gives, including the
value $+\infty$,
\begin{equation}
 a_2^T\|B_z\xi\|^2
 =c_i^{2\vartheta}\int_{B_S^{\mathrm{ch}}}
 K_S(y(Z),z)^2
 \big\|b_{y(Z)}^{-\vartheta}(VI\xi)(Z)\big\|^2\,dZ.
 \label{eq:locality-graph-formula}
\end{equation}
Here $B_S^{\mathrm{ch}}$ denotes the full chamber variable occurring
in $V$; its indicated vertex marginal is the ordinary probability
on $B_S$.  The left-hand side follows by integrating $v$ and using
Lemma~\ref{lem:locality-kernel}.  Integrating the right-hand side in
$z$ replaces the kernel square by the constant $a_2^S$.  This
identifies the graph norm with a closed weighted norm of $VI\xi$,
and proves completeness.  Boundedness of $K_S$ gives the first
estimate.  For the second, the triangle inequality implies that at
least one of $d(y,z),d(y,w)$ is at least $d(z,w)/2$.  The sum of the
two kernel squares therefore dominates $c\,d(z,w)^8$, uniformly in
$y$, giving Equation~\eqref{eq:input-24}.

If $\sum_j a_j n_S(z_j)=0$, then
$\sum_j a_jK_S(y,z_j)=0$ for every $y$.  All the source operators in
this finite linear combination are defined on $VI\mathfrak D$ by
Equation~\eqref{eq:locality-graph-formula}.  Intertwining shows that
$\sum_j a_jA_{z_j}(v)\xi=0$ for almost every $v$.  Continuity in $v$
on this domain extends the equality to every $v$.  Moreover,
\[
 B_z\xi=(a_1^T)^{-1}\int_{B_T}A_z(v)\xi\,dv,
 \qquad \xi\in\mathfrak D,
\]
so the same relation holds for $B_z$.  Linearity in $n_T(v)$ follows
directly from the bilinear kernel representation.
\end{proof}

\begin{lemma}[Density, smooth stability, and cores]
\label{lem:locality-cores}
The domain $\mathfrak D$ is dense in $\mathcal H$.  For each compact
subset $C\subset B_T\times\mathbb R$ there are $r\in\mathbb N$ and
$C_C<\infty$ such that
\[
 \|f(Q_z)\xi\|_{\mathfrak D}
 \le C_C\|f\|_{C^r}\|\xi\|_{\mathfrak D},
 \qquad \supp f\subset C.
\]
There are $\chi_l\in C_c^\infty(B_T\times\mathbb R)$ such that
$\chi_l(Q_z)\xi\to\xi$ in $\mathfrak D$ for every
$\xi\in\mathfrak D$.  Finally, $\mathfrak D$ is a core for every
$B_z$ and $A_z(v)$.
\end{lemma}

\begin{proof}
\emph{Density.}
Let $M_l$ be multiplication on the source chamber-variable space by
$b_y^\vartheta1_{\{b_y\le l\}}$, and put
\[
 T_l=I^*V^*M_lI_S.
\]
These are bounded operators.  The constant-input formula in
Equation~\eqref{eq:input-22} and the strong commutation of
$D_Z$ with its own-coordinate height give
\[
 \langle\eta,T_l\eta\rangle
 =c_0^{-1/2}\int
 \left\langle D_Z^{-1/2}\eta,
 b_y^\vartheta1_{\{b_y\le l\}}\eta\right\rangle\,dZ\ge0.
\]
All pairings are integrable by the isometry and boundedness of $M_l$.
For nonzero $\eta$, they cannot vanish for every $l$: in the common
$D_Z,b_y$ calculus the union of the cuts $b_y\le l$ is the identity,
and both positive operators have zero kernel.  Thus the positive
operators $T_l$ have zero common kernel and their ranges are total.

For $\zeta=V^*M_lI_S\eta$, the reducing-range intertwining gives
\[
 A_z\zeta
 =V^*C_z M_lI_S\eta
 =c_i^\vartheta V^*
   \big(K_S(y,z)1_{\{b_y\le l\}}I_S\eta\big).
\]
Its norm is bounded uniformly in $z$.  To estimate $B_zI^*\zeta$,
test against $\eta'\in\operatorname{dom} B_z$ and write
\[
 \langle I^*\zeta,B_z\eta'\rangle
 =\int_{B_T}
 \left\langle A_z(v)\zeta(v),
 K_T(v,P_z)^{-1}\eta'\right\rangle\,dv.
\]
This identity is first made with joint spectral truncations of $Q_z$.
The inverse-square bound in Lemma~\ref{lem:locality-kernel} permits
their removal by Cauchy--Schwarz.  It proves that $I^*\zeta\in\operatorname{dom} B_z$
with a uniform bound.  Integrating $z$ proves $\ran T_l\subset\mathfrak D$,
and therefore density.

\emph{Smooth stability and graph cutoffs.}
The basic stability estimate concerns a bounded unary average
\[
 R_F=\int_{B_T}F(v,A_z(v))\,dv.
\]
For $\xi\in\mathfrak D$ and $\eta'\in\operatorname{dom} B_w$, move $A_w(v)$
through $F(v,A_z(v))$, using
Lemma~\ref{lem:locality-graph-identities}, and pair against
$K_T(v,P_w)^{-1}\eta'$.  Cauchy--Schwarz gives
\begin{equation}
 \|B_w R_F\xi\|
 \le C\|F\|_\infty\|B_w\xi\|.
 \label{eq:locality-unary-bound}
\end{equation}
The adjoint domain criterion used above justifies this computation
before the left-hand side is known to be finite.  Integrating $w$
gives a bound on $\mathfrak D$.  Products of unary averages satisfy
the iterated bound.

Choose $v_1,\ldots,v_d$ such that the $n_T(v_j)$ form a basis of
$E_T$.  Allow each $v_j$ to vary in a sufficiently small fixed
neighborhood.  The resulting inverse frame matrices are uniformly
bounded.  The evaluations
\[
 \big(K_T(v_j,u)e^{\vartheta s}\big)_{j=1}^d
\]
then recover $e^{\vartheta s}n_T(u)$ linearly.  A compactly supported
smooth function of $(u,s)$ extends smoothly to a compact set in the
ambient vector space away from zero.  Compose such an extension with
the linear recovery, and integrate the independent $v_j$ over their
neighborhoods.  On fixed compact coordinate boxes, its Fourier
expansion in the separate variables has absolutely summable
coefficients after finitely many derivatives.  Each separated term
is a product of unary averages.  Equation~\eqref{eq:locality-unary-bound}
therefore proves the stated smooth stability estimate.

For completeness, compact graph cutoffs can be constructed by this
same procedure, without assuming a joint height law at different
labels.  Cut each of the basis evaluations smoothly from above at
$l$, and subtract the piece on which all evaluations are smaller
than $1/l$.  Uniform frame inverses imply that the resulting
function of $e^{\vartheta s}n_T(u)$ is supported in a bounded radial
interval bounded away from zero.  It is consequently a compact
smooth function of $(u,s)$.  Each upper unary cutoff tends strongly
to the identity.  Each simultaneous-small piece tends strongly to
zero.  The same convergences hold in the graph norm: in the proof of
Equation~\eqref{eq:locality-unary-bound}, apply dominated convergence
to the vector $A_w(v)\xi$, then integrate $w$.  The only possible
zero of an evaluation at a finite positive height is $P_z= v$;
$P_z(\{v\})=0$ for almost every ordinary $v$, as follows by integrating
against a faithful normal scalarization of $P_z$ and using that the
ordinary measure is nonatomic.

\emph{Individual cores.}
Let $L$ be $B_z$ or $A_z(v)$.  Cut a
vector in $\operatorname{dom} L$ by compact smooth functions of its own $Q_z$
variable that tend to one and have bounded height support.  These
cuts converge in the graph norm of $L$.  For each fixed cutoff,
approximate the uncut vector in Hilbert norm by vectors of the dense
domain $\mathfrak D$, and apply that cutoff.  Smooth stability keeps
the resulting vectors in $\mathfrak D$, while boundedness of $L$
on the cutoff gives convergence in its graph norm.  A diagonal
choice proves that $\mathfrak D$ is a core.
\end{proof}

\subsection{Exact locality and the exclusion of derivatives}

\begin{lemma}[Real powers on transverse patches]
\label{lem:locality-real-substitution}
Fix $z\ne w$.  On compact transverse angular patches, with the
$Q_z$ height restricted to a compact interval, the ordered
replacement
\[
 B_w\quad\longmapsto\quad
 c_i^\vartheta K_S(w,z)K_T(u',u)^{-1}B_z^{-1}
\]
is valid, where the angular variables occur in the order
$(u',u)=(P_w,P_z)$.  The adjoint ordered replacement is valid as
well.  After division by $K_S(w,z)$ the norms of the replaced
operators are bounded independently of $w$; the fixed angular and
height patches determine the bound.
\end{lemma}

\begin{proof}
To specify the ordered assertion, let $f\in C_c^\infty(B_T\times\mathbb R)$,
let $\rho$ equal one on its angular support, and let $\chi'$ be supported
in a patch transverse to $\supp\rho$. For a smooth kernel on this
transverse rectangle, write
\[
 \mathcal M_{w,z}(F)
   =\int F(u',u)\,P_w(du')P_z(du).
\]
This denotes the operator-norm-convergent sum obtained from a summable
separated smooth expansion of $F$; it does not introduce joint spectral
calculus for $P_w,P_z$. We will prove the bounded operator identity
\begin{equation}
 B_wP_w(\chi')f(Q_z)
 =c_i^\vartheta K_S(w,z)\,
   \mathcal M_{w,z}\!\left(
        \frac{\chi'(u')\rho(u)}{K_T(u',u)}\right)
   B_z^{-1}f(Q_z).
 \label{eq:locality-real-block}
\end{equation}
In particular the entire range on the left belongs to the domain of
$B_w$. Initially insert a fixed rightmost smooth angular-height test
$f_0(Q_z)$ equal to one on $\supp f$, with a slightly larger transverse
angular support and a fixed compact height interval. Equation~\eqref{eq:input-19}
then gives, for a Schwartz function $a$ and by Fourier inversion,
\[
 \begin{aligned}
 &a(\log b_w)P_w(\chi')f_0(Q_z)\\
 &\quad=\int \chi'(u')\rho_0(u)
 a\!\left(\log\frac{c_i k_S(w,z)}{k_T(u',u)}-s\right)
 \,P_w(du')Q_z(du,ds)f_0(Q_z),
 \end{aligned}
\]
where $\rho_0$ is a fixed enlarged angular cutoff. In this formula the
height on the right belongs to the single $Q_z$ calculus, and the
remaining ordered kernel is defined by the same separated expansions.
Every fixed number of derivatives of $k_T(u',u)^{-\mathrm it}$ on the
transverse rectangle is bounded polynomially in $|t|$. The Schwartz
Fourier transform therefore makes the integral and all required
separated expansions absolutely summable.

The substituted real logarithm belongs to an interval
$\log(c_i k_S(w,z))+J$, where the compact interval $J$ depends only on
the fixed supports. Every smooth $a$ supported outside this interval
annihilates the left block. By spectral regularity this proves that
$P_w(\chi')f_0(Q_z)$ has its $\log b_w$ spectrum in that interval.
This is a bounded-operator support statement on the whole input range.
It applies in particular after inserting any additional bounded
oscillatory multiplier on that range, with a bound independent of its
derivatives. We may now insert $a(s)=e^{\vartheta s}$ on the indicated
interval using spectral truncations. This proves
Equation~\eqref{eq:locality-real-block}. After division by $K_S(w,z)$,
the remaining angular kernel is fixed and smooth and $B_z^{-1}$ is
bounded on the fixed input height interval. The operator norm bound is
therefore uniform in $w\ne z$. Taking adjoints gives the corresponding
ordered substitution with a fixed output test when the output is the
transverse side.
\end{proof}

\begin{lemma}[Angle and height locality]
\label{lem:locality-support}
Let $f,g$ be compact smooth functions of $Q_z$ whose angular
supports are disjoint.  For every $w\in B_S$, $v\in B_T$, and
$\xi,\eta\in\mathfrak D$,
\begin{equation}
 \langle g(Q_z)\eta,A_w(v)f(Q_z)\xi\rangle=0.
 \label{eq:input-25}
\end{equation}
For almost every $v$, the vector $A_w(v)f(Q_z)\xi$ is supported, in
the full angle-height representation of $Q_z$, in $\supp f$.
The same support assertion holds for $B_wf(Q_z)\xi$.
\end{lemma}

\begin{proof}
Partition the middle angular variable $P_w$ into finitely many
smooth pieces so that each is transverse either to the angular
support of $f$ or to that of $g$.  Such a partition exists because
those two supports are disjoint and, in rank one, failure of
transversality is equality.  For $w\ne z$, apply
Lemma~\ref{lem:locality-real-substitution} on the transverse side of
each piece.  After dividing the resulting scalar matrix element by
$K_S(w,z)$, all remaining factors have fixed smooth separated
expansions and bounded $Q_z$ heights.  As $w\to z$, strong
continuity of the angular PVM makes that divided matrix element
tend to zero: at $w=z$ the angular slots are in the same calculus,
and the disjoint outer supports annihilate.

On the other hand, the original scalar matrix element is linear in
$n_S(w)$ by Lemmas~\ref{lem:locality-graph-identities} and
\ref{lem:locality-cores}.  Use a unipotent chart with $z$ at its
origin.  After multiplication by a fixed smooth nonzero
denominator, this is a polynomial of weighted degree at most four.
Along the graded dilations of any fixed chart point, the preceding
limit says that it is $o(t^4)$.  Every homogeneous part of degrees
zero through four must therefore vanish.  The polynomial is zero
throughout the chart, and continuity extends this to the whole
boundary.  This proves Equation~\eqref{eq:input-25}, including
$w=z$.

For fixed $v$, $A_w(v)$ strongly commutes with $A_z(v)$.  The latter
is the multiplier
\[
 m_v(u,s)=K_T(v,u)e^{\vartheta s}
\]
in the $Q_z$ representation.  At every angle $u\ne v$, two different
heights can be separated by disjoint intervals of this multiplier,
after shrinking their angular neighborhoods.  Combine this
separation with Equation~\eqref{eq:input-25} and a finite smooth
partition on compact sets.  Strong commutation with the spectral
cuts of $m_v$ eliminates the blocks with separated heights, while
Equation~\eqref{eq:input-25} eliminates the blocks with separated
angles.  The exceptional output projection $P_z(\{v\})$ vanishes
for almost every $v$.  Exhaustion of the complement of $\supp f$
proves full angle-height locality.  Finally integrate the operators
$A_w(v)$ on $\mathfrak D$ and use
$B_w=(a_1^T)^{-1}\int A_w(v)\,dv$.
\end{proof}

The preceding locality statement still allows differential
operators.  The next two arguments eliminate them and also explain
why no choice of a self-adjoint extension is hidden in the passage
to strong commutation.

\begin{lemma}[Finite jets from graph locality]
\label{lem:locality-jets}
Fix $\xi\in\mathfrak D$ and a relatively compact coordinate chart
in $B_T\times\mathbb R$.  In the spectral representation of $Q_z$,
write its variable as $q$.  For $L=B_w$ and for a countable dense
set of the operators $L=A_w(v)$ for which
Lemma~\ref{lem:locality-support} holds, there is a common finite
integer $m$ such that, on a smaller chart,
\begin{equation}
 Lf(Q_z)\xi
 =\sum_{|\alpha|\le m}\eta_{L,\alpha}(q)
                          \partial^\alpha f(q),
 \qquad f\in C_c^\infty,
 \label{eq:locality-jet-expansion}
\end{equation}
where the coefficient vectors are measurable and locally square
integrable in the spectral measure.  The same assertion holds for
any finite family of operators with these graph bounds and locality.
\end{lemma}

\begin{proof}
Let $\mathcal Q=B_T\times\mathbb R$, and realize the spectral representation as
\[
 \mathcal H=\int_{\mathcal Q}^{\oplus}\mathcal H_q\,d\nu(q),
 \qquad (f(Q_z)\eta)(q)=f(q)\eta(q).
\]
Here $\nu$ is a finite control measure obtained from a faithful normal
scalarization; it is not assumed to be related to smooth volume on $\mathcal Q$.
Fix a compact coordinate box and a smooth cutoff equal to one on a
smaller box. Lemma~\ref{lem:locality-cores}, followed by the bound
$\|L\eta\|\le C_L\|\eta\|_{\mathfrak D}$, gives
\[
 \|Lf(Q_z)\xi\|\le C_{L,\xi}\|f\|_{C^r}
\]
for tests in the box. Take a Sobolev order $a>r+\dim(\mathcal Q)/2$ with
respect to ordinary coordinate volume. Multiplication by the fixed cutoff
makes this a bounded map from $H^a$ on a containing coordinate torus to
$\mathcal H$. Increase to an integer $b$ with $b-a>\dim(\mathcal Q)/2$.
The inclusion $H^b\hookrightarrow H^a$ is Hilbert--Schmidt: on the
Fourier basis its squared singular values have convergent sum
$\sum_{k\in\mathbb Z^{\dim\mathcal Q}}(1+|k|^2)^{-(b-a)}$.
Consequently the map $T_L:f\mapsto Lf(Q_z)\xi$, including the cutoff,
is Hilbert--Schmidt from $H^b$ to $\mathcal H$.

If $(e_j)$ is an orthonormal basis of $H^b$, choose measurable versions
$t_j(q)=(T_Le_j)(q)$. Tonelli gives
\[
 \int_{\mathcal Q}\sum_j\|t_j(q)\|^2\,d\nu(q)
     =\sum_j\|T_Le_j\|^2<\infty.
\]
Outside one null set, $f\mapsto\sum_j\langle e_j,f\rangle_{H^b}t_j(q)$
is therefore a continuous $\mathcal H_q$-valued distribution, and it
represents $(T_Lf)(q)$ for all smooth tests in a countable dense family.
Continuity extends the equality to all tests. Choose that family to
include dense tests in every box of a countable coordinate basis.
Locality makes the distribution vanish on each such box not containing
$q$, so its support is contained in $\{q\}$. A $C^b$ seminorm bounds
its $H^b$ norm on the fixed box. Thus its distribution order is at most
the same finite integer $b$, independently of $q$.

This is the distributional finite-jet mechanism underlying Peetre's
locality principle \cite{Peetre1959,Peetre1960}; compare the scalar
point-support theorem \cite[Theorems~2.3.3--2.3.4]{Hormander2003}.
The preceding construction uses the actual control measure and does
not require it to be absolutely continuous. For completeness, a
point-supported vector distribution of order at most
$b$ annihilates any smooth function whose derivatives through order $b$
vanish at the point. Multiply the Taylor remainder by a cutoff of radius
$\epsilon$ about that point. All its derivatives through order $b$ tend
to zero uniformly as $\epsilon\downarrow0$, whereas support at the point
leaves its distribution value unchanged. This proves the finite-jet
formula~\eqref{eq:locality-jet-expansion} with $m=b$. On a smaller box,
cutoff monomials $x^\alpha/\alpha!$, $|\alpha|\le b$, recover the
coefficients by a triangular system with bounded polynomial entries in
$q$. Each coefficient is therefore a finite combination of the vectors
$L f_\alpha(Q_z)\xi$ with bounded measurable coefficients. In particular
it is measurable and locally square integrable for $\nu$.

A common finite order for the whole family is essential here. Choose
finitely many source labels whose lifts span $E_S$, and choose a target
lift basis from the common full-measure set on which locality holds for
those source labels. All $A_w(v)$ on $\mathfrak D$ are linear combinations
of this finite array, because their dependence on $n_S(w)$ and $n_T(v)$
is bilinear. All $B_w$ are combinations of the same array by their
ordinary-$v$ averages. Apply the preceding construction to this finite
array and take the largest Sobolev and jet order. Its formulas and
locality then extend by finite linearity to every operator under
consideration. This takes no supremum over an uncontrolled countable
family of differential orders.
\end{proof}

\begin{lemma}[Vanishing of positive differential order]
\label{lem:locality-symbols}
Every expansion in Equation~\eqref{eq:locality-jet-expansion} has
order zero.  Consequently the $Q_z$ and $Q_w$ calculi strongly
commute for all source labels $z,w$.
\end{lemma}

\begin{proof}
Use a common upper order $m$ and suppose $m>0$.  The highest
homogeneous symbol is the vector polynomial
\[
 \sigma_L(q,\zeta)
   =\sum_{|\alpha|=m}\eta_{L,\alpha}(q)\zeta^\alpha.
\]
It is obtained in Hilbert norm by applying $L$ to
$e^{\mathrm i t\psi}\chi$, removing the phase, and dividing by
$(\mathrm i t)^m$, where $\psi$ is real and smooth and $\chi$ is a
compactly supported smooth cutoff.  Terms with fewer derivatives
vanish in this limit.

For $w\ne z$, take the angular support of $\chi$ in a small ball
about $u_0$.  By Lemma~\ref{lem:locality-real-substitution}, the
$P_w$ component outside a slightly larger ball has bounded norm
after applying $B_w$ to
$e^{\mathrm i t\psi}\chi(Q_z)\xi$, independently of $t$.  On the
remaining component,
$A_w(v)-K_T(v,u_0)B_w$ costs at most the variation of
$K_T(v,\cdot)$ on that enlarged ball times the $B_w$ output norm.
After dividing by $t^m$ and taking the limit, this gives
\[
 \|\big(\sigma_{A_w(v)}(q,d\psi)
       -K_T(v,u_0)\sigma_{B_w}(q,d\psi)\big)\chi\|_2
 \le \operatorname{var}(K_T)
       \|\sigma_{B_w}(q,d\psi)\chi\|_2.
\]
Use arbitrary inner cutoffs and rational linear phases in the chart.
Regularity of the spectral control measure turns the norm inequality
for every such cutoff into the corresponding almost-everywhere pointwise
inequality on the ball. Take a countable family of shrinking balls; their
variations tend to zero and their centers tend to the tested angle.
This gives the following polynomial identity without requiring a
Lebesgue differentiation theorem for the spectral measure:
\begin{equation}
 \sigma_{A_w(v)}(q,\zeta)
       =K_T(v,u(q))\sigma_{B_w}(q,\zeta)
 \label{eq:locality-symbol-comparison}
\end{equation}
almost everywhere, as a polynomial identity.  For $w=z$ the
operators already are multiplication operators, so their
positive-order symbols are zero.

Strong commutation of $A_w(v)$ with $A_z(v)$ is now used a second
time.  On the support of a compact test, the smooth multiplier
$m_v(q)=K_T(v,u(q))e^{\vartheta s(q)}$ is bounded.  Locality places
the output on this same compact support.  Thus both sides of the
commutation identity are defined, and substitution of
Equation~\eqref{eq:locality-jet-expansion} into that identity is
legitimate.  The coefficients of derivatives of order $m-1$ in the
commutator give
\[
 D_{d m_v(q)}\sigma_{A_w(v)}(q,\zeta)=0,
\]
where $D_{d m_v(q)}$ denotes differentiation in the covariable in
the direction of the covector $d m_v(q)$.  All lower-order terms
cancel or have smaller derivative order.  For $v\ne u(q)$, the
factor $K_T(v,u(q))$ in
Equation~\eqref{eq:locality-symbol-comparison} is positive.  Hence
\[
 D_{d m_v(q)}\sigma_{B_w}(q,\zeta)=0.
\]
The functions $m_v(q)$ are the pairings of the embedded cone point
$e^{\vartheta s}n_T(u)$ with the spanning vectors $n_T(v)$.
Their differentials, with $v\ne u$, span the cotangent space of the
angle-height manifold.  Therefore $\sigma_{B_w}$ is independent of
the covariable.  Homogeneity of positive degree makes it zero;
Equation~\eqref{eq:locality-symbol-comparison} does the same for
the other symbols.  Decrease $m$ and repeat until only order zero
remains.

The coefficient formula concerns the map $f\mapsto Lf(Q_z)\xi$ for a
fixed $\xi\in\mathfrak D$. Once its order is zero, it gives
\[
 L(F\chi)(Q_z)\xi=F(Q_z)L\chi(Q_z)\xi
\]
for compact smooth tests supported in one coordinate chart. A finite
partition of unity gives the same equality for arbitrary compact smooth
$F,\chi$. Take $\chi$ from the graph cutoffs of
Lemma~\ref{lem:locality-cores}. Then $\chi(Q_z)\xi\to\xi$ in the
graph norm of $L$, so closedness gives
\[
 LF(Q_z)\xi=F(Q_z)L\xi,\qquad \xi\in\mathfrak D.
\]
Since $\mathfrak D$ is a core for the self-adjoint operator $L$, another
graph approximation extends this equality to every $\xi\in\operatorname{dom}L$.
Applying it to $(L-\mathrm i)^{-1}\eta$ shows that $F(Q_z)$ commutes with
the resolvent. Spectral regularity and bounded strong approximation by
compact smooth functions now give commutation with every projection of $Q_z$. This proves strong commutation for $L=B_w,A_w(v)$. A finite basis
of the evaluations $A_w(v)$ recovers the cone vector $B_wn_T(P_w)$, hence
both $P_w$ and $\log b_w$. The complete $Q_w$ calculus therefore commutes
with $Q_z$.
\end{proof}

\subsection{The resulting cone-map law}

\begin{lemma}[Smooth angular bijections]
\label{lem:locality-cone-maps}
The common spectral law of the $Q_z$ is a law of evaluations
\[
 z\longmapsto (\phi(z),\log h(z)),
\]
where $h$ is positive and smooth, $\phi:B_S\to B_T$ is a smooth
diffeomorphism, and
\begin{equation}
 U n_S(z)=c_i^{-\vartheta/2}h(z)^\vartheta n_T(\phi(z))
 \label{eq:locality-cone-isometry}
\end{equation}
for a linear pairing isometry $U:E_S\to E_T$ carrying the positive
cone onto the positive cone.  The space $\mathcal X$ of these
linear isometries is a smooth Lie coset, and the spectral law may
be regarded as a PVM on $\mathcal X$.
\end{lemma}

\begin{proof}
Strong spectral continuity in $z$, supplied by compact source transport,
gives versions measurable jointly in the spectral variable and the ordinary source label.
For example, use a faithful normal scalarization of the common
abelian algebra, approximate in measure on a countable dense label
set, and take a fast almost-everywhere convergent subsequence of
the resulting simple fields.  For generic distinct $z,w$, their
images cannot agree on a nonzero spectral projection $e$. Otherwise a
nonzero vector $\eta\in e\mathcal H$ satisfies
$P_z(f)\eta=P_w(f)\eta$ for every continuous $f$. On finite smooth sums
$b(u)=\sum_j f_j(u)\xi_j$, the functional
\[
 b\longmapsto\left\langle\eta,\sum_jP_z(f_j)\xi_j\right\rangle
     =\left\langle\eta,\sum_jP_w(f_j)\xi_j\right\rangle
\]
is nonzero, by taking $b$ constant equal to $\eta$, and is bounded by
both sharp form norms in Equation~\eqref{eq:input-3}. This contradicts
the disjointness proved in Section~\ref{sec:transfer}. Exhaust transverse
angular rectangles in Equation~\eqref{eq:input-19} and take a common
conull set for rational frequencies. The resulting equality of scalar
unitary characters gives, almost surely and for almost every ordinary pair,
\begin{equation}
 h(z)^\vartheta h(w)^\vartheta
 K_T(\phi(z),\phi(w))=c_i^\vartheta K_S(z,w).
 \label{eq:locality-classical-pairing}
\end{equation}

Fubini permits two fixed anchors $z_0,z_1$ with distinct images almost
surely, and a fixed generic countable dense set of source labels on which
all pair identities and all identities with the anchors hold almost
surely. Such a dense set can be selected by choosing one generic point
in each member of a countable boundary basis. Fix a realization in this
common conull set and put $r(z)=h(z)^\vartheta$. Adding the two anchor
identities gives
\[
 r(z)\big(r(z_0)K_T(\phi(z),\phi(z_0))
          +r(z_1)K_T(\phi(z),\phi(z_1))\big)
 =c_i^\vartheta\big(K_S(z,z_0)+K_S(z,z_1)\big).
\]
The sum of the target kernels against their distinct images has a positive
minimum on the compact target boundary.  The two identities and
boundedness of the source kernels give an upper bound for
$h(z)^\vartheta$.  Conversely, the sum of the source kernels
against the two distinct anchors has a positive minimum, and
boundedness of the target kernels gives a positive lower bound for
$h(z)^\vartheta$.

On the chosen countable dense set, Lemma~\ref{lem:locality-kernel} and
the two bounds on $h$ give
\[
 c\,d_S(z,w)\le d_T(\phi(z),\phi(w))\le C\,d_S(z,w),
\]
where the constants may depend on this realization. Thus $\phi$ and
its inverse on its image are uniformly continuous.  Metric completion gives a
continuous injective map on the whole boundary.  The displayed sum of anchor pairings has a strictly positive
denominator everywhere and recovers a continuous positive extension
of the weight.  Pair identities with almost every further source point
show that these extensions give versions of the original fields.

A continuous injection between spheres implies that the source
dimension is at most the target dimension.  This conclusion holds
at every isolated type.  The type correspondence is a permutation
of a finite list, so the inequalities around each permutation
cycle are equalities.  Brouwer's invariance of domain (see
\cite[Theorem~2 and Corollary~3]{Tao2011Invariance} for the openness
and dimension statements) now makes $\phi(B_S)$
open in $B_T$; it is also compact and nonempty, hence all of $B_T$.

Define $U$ on the spanning set $n_S(B_S)$ by
Equation~\eqref{eq:locality-cone-isometry}.  It is well-defined
linearly.  Indeed a linear relation among the source lifts pairs
to zero against every source lift.  Equation~\eqref{eq:locality-classical-pairing}
makes its proposed image pair to zero against every target lift,
because $\phi$ is onto.  Nondegeneracy and spanning force that
image relation to be zero.  The same equation proves that $U$
preserves the pairing, and its image spans $E_T$, so it is an
isomorphism.  It carries the positive radial cone onto the positive
radial cone.  The smooth embedded cone descriptions now imply that
$\phi$, $h$, and the angular inverse are smooth.

Choose one such $U_0$.  Then $\mathcal X$ is the coset of the
subgroup of the target orthogonal group preserving its positive
cone, multiplied by $U_0$.  This subgroup is closed: preservation
of the cone with its vertex follows by limits, and the same
condition on inverse matrices gives equality rather than mere
inclusion.  The Closed Subgroup Theorem makes it a Lie group and
$\mathcal X$ a smooth embedded coset.  Choose a lift basis from the fixed generic dense source set. Its finitely
many evaluation vectors determine $U$ by a fixed linear inversion. Thus
the passage from the original measurable map version to $U$ is Borel,
with no measurable choice of a continuously varying boundary map
required.  Conversely evaluation and
angular inversion on this cone are continuous.  The law on
$\mathcal X$ recovers every fixed label by continuity in measure.
\end{proof}

\subsection{Two normal copies sharing the trunk}

Use superscripts $1,2$ for the two measured copies. The same-label
statement following Equation~\eqref{eq:input-20} in
Section~\ref{sec:heights} gives strong commutation of $Q_z^1$ with $Q_z^2$.
It includes the mixed angle--height relations and the height groups; for
the latter, interchange of the two normal-copy slots makes the possible
scalar bicharacter vanish. Each copy separately has the $\mathcal X$
representation of Lemma~\ref{lem:locality-cone-maps}.

To compare these two representations, we need a domain on which compact
smooth functions of the first-copy matrix preserve the second-copy
heights. We obtain it by averaging the available same-label commutation.
The inverse-height estimate below plays the role of the inverse-kernel
estimate used to construct the first graph domain. Its integrability
again comes from fourth-order boundary geometry and $Q>8$.

In the second calculus put
\[
 d_z=(b_z^2)^\vartheta,
 \qquad H_2=\left(\int_{B_S}d_z^2\,dz\right)^{1/2},
 \qquad \mathfrak B=\operatorname{dom}H_2.
\]
The superscript on $b_z^2$ denotes the copy, not a square.

\begin{lemma}[An inverse-height estimate and first-copy graph cuts]
\label{lem:locality-two-copy-domain}
In the second classical calculus,
\begin{equation}
 d_z\le C H_2,\qquad
 d_z+d_w\ge c_2d(z,w)^4H_2,
 \qquad
 \int_{B_S}(H_2/d_z)^2\,dz\le C.
 \label{eq:locality-inverse-height}
\end{equation}
The domain $\mathfrak B$ is dense and is a core for all $d_z$ and
$A_z^2(v)$. Compact smooth functions of the first-copy $\mathcal X$
preserve $\mathfrak B$, with a finite differentiability bound on each
compact set. There are compact smooth first-copy cutoffs converging to
the identity in the graph norm of $\mathfrak B$.
\end{lemma}

\begin{proof}
\emph{The inverse-height estimate and the second-copy cores.}
The first two inequalities follow from Equation~\eqref{eq:input-24},
using commutation in the second copy and continuity in the source label.
In a realization of this calculus write $r_z=d_z/H_2$. If
$r_z<\epsilon$ and $r_w<\epsilon$, the second inequality gives
$d(z,w)\le C\epsilon^{1/4}$. The sublevel set thus has ordinary measure
at most $C\epsilon^{Q/4}$. Dyadic decomposition gives
\[
 \int r_z^{-2}\,dz
 \le 1+4\sum_{j\ge0}2^{2j}|\{z:r_z\le2^{-j}\}|
 \le 1+C\sum_{j\ge0}2^{(2-Q/4)j}<\infty.
\]
The constants are uniform in the realization, proving the third estimate.
Each second-copy cone map has bounded positive weights on the compact
source boundary. Thus $H_2$ is finite and positive in every realization,
and its bounded spectral cuts prove density of $\mathfrak B$. These cuts
commute with each $d_z,A_z^2(v)$ and converge in their individual graph
norms. Consequently $\mathfrak B$ is a core for each of them.

\emph{Unary averages and smooth matrix functions.}
For a bounded measurable unary field $F$, put
\[
 S_F=\int_{B_S}F(z,Q_z^1)\,dz.
\]
For $\xi\in\mathfrak B$ and initially $\eta\in\operatorname{dom}H_2$,
same-label commutation gives
\[
 \langle S_F\xi,H_2\eta\rangle
 =\int_{B_S}
 \left\langle F(z,Q_z^1)d_z\xi,(H_2/d_z)\eta\right\rangle\,dz.
\]
Spectral truncation justifies this equality. The inverse-height bound and
Cauchy--Schwarz bound its right-hand side by
\[
 C\left(\int_{B_S}\|F(z,Q_z^1)d_z\xi\|^2\,dz\right)^{1/2}\|\eta\|.
\]
The adjoint-domain criterion therefore gives $S_F\xi\in\mathfrak B$ and
\begin{equation}
 \|H_2S_F\xi\|
 \le C\left(\int_{B_S}\|F(z,Q_z^1)d_z\xi\|^2\,dz\right)^{1/2}
 \le C\|F\|_\infty\|H_2\xi\|.
 \label{eq:locality-two-copy-unary}
\end{equation}
Only the same-label commutation has been used.

Choose source-label neighborhoods $O_1,\ldots,O_d$ such that
$n_S(z_1),\ldots,n_S(z_d)$ form a basis whenever $z_j\in O_j$,
with uniformly bounded inverse basis matrices. Here $d=\dim E_S$.
Fix smooth probability densities $\rho_j$ supported in these neighborhoods.
In fixed bases of $E_S,E_T$, the evaluation vectors recover the matrix by
\[
 U=[U n_S(z_1),\ldots,U n_S(z_d)]
       [n_S(z_1),\ldots,n_S(z_d)]^{-1}.
\]
Extend a compact smooth function on $\mathcal X$ to the ambient matrix
space and compose it with this formula. The result is smooth in the
separate pairs $(z_j,U n_S(z_j))$ on fixed buffered boxes. It has uniform
compact support in the evaluation variables because the basis matrices
and their inverses are uniformly bounded. A summable separated expansion,
followed by independent integration against the $\rho_j$, expresses its
first-copy calculus as a sum of products of the unary averages above.
Equation~\eqref{eq:locality-two-copy-unary} proves the required finite
smooth-seminorm bound on $\mathfrak B$.

\emph{Compact cutoffs with a uniform graph bound.}
Choose $\chi\in C_c^\infty(\mathbb R)$ with $0\le\chi\le1$ and
$\chi=1$ near $[0,1]$. For $l\ge1$ define first-copy multipliers
\[
 t_{j,l}(U)=\int_{B_S}\rho_j(z)
   \chi\!\left(\|U n_S(z)\|^2/l^2\right)\,dz,
 \qquad \chi_l(U)=\prod_{j=1}^d t_{j,l}(U).
\]
Let $T_{j,l}$ be the first-copy spectral multiplier with scalar function
$t_{j,l}$. Each integrand is a bounded function of $Q_z^1$, since
$U n_S(z)=c_i^{-\vartheta/2}(b_z^1)^\vartheta n_T(P_z^1)$.
The graph bounds for $T_{j,l}$ depend on the fixed $\rho_j$ and $\chi$,
and are independent of $l$. Applying the first inequality in
Equation~\eqref{eq:locality-two-copy-unary} to the field
\[
 \rho_j(z)\left[
 \chi\!\left(\|U n_S(z)\|^2/l^2\right)-1\right]
\]
shows, by dominated convergence on the vectors $d_z\xi$, that
$T_{j,l}\xi\to\xi$ in the graph norm of $H_2$. The finite product has
the same convergence, by telescoping and the uniform graph bounds.

Finally, $\chi_l$ is a compact smooth function on $\mathcal X$.
Indeed, if its product integral is nonzero, some tuple of labels in the
fixed neighborhoods has every evaluation norm bounded by a constant
times $l$. The recovery formula then bounds $\|U\|$ by a constant times
$l$. Bounded subsets of the pairing-isometry space have bounded inverses,
since the inverse is obtained from the transpose using the two fixed
nondegenerate pairing matrices. The cone-preserving condition is closed,
so these closed bounded matrix windows are compact in $\mathcal X$.
This proves the asserted compact graph cutoffs.
\end{proof}

\begin{lemma}[Cross-copy strong commutation]
\label{lem:locality-two-copy-commutation}
The two normal $\mathcal X$ measurements strongly commute.
\end{lemma}

\begin{proof}
We first establish locality in the first-copy $\mathcal X$
representation for $L=d_z,A_z^2(v)$ on $\mathfrak B$.  If two
matrices $U,U'$ are distinct, some source evaluation separates
them.  After shrinking matrix neighborhoods and a neighborhood of
that label, their entire evaluation images are disjoint for every
label in this open set.  Between smooth tests supported in these
two matrix neighborhoods, the blocks of $d_z$ and $A_z^2(v)$ vanish
for those labels by same-label commutation.  The scalar blocks are
linear in $n_S(z)$ on $\mathfrak B$.  A linear functional vanishing
on the lifts of a nonempty open boundary patch vanishes on their
whole span: in a unipotent chart it is a polynomial divided by a
nonzero denominator, and polynomial uniqueness applies.  Hence
these blocks vanish for every $z$.  Covering any pair of disjoint
compact matrix supports by finitely many such separating
neighborhoods proves locality.

Apply the proof of Lemma~\ref{lem:locality-jets} on a compact chart
of $\mathcal X$, now using the graph estimates of
Lemma~\ref{lem:locality-two-copy-domain}.  Finite-dimensional
linearity again gives a common finite upper order for a spanning
family and thus for a countable dense list of labels and
evaluations. We fix these countable lists, their rational coordinate
phases, and a countable chart exhaustion at the outset; all pointwise
polynomial assertions can then be made on one conull set. Suppose the
current common highest order is positive.
Same-label commutation says that the principal symbol of $d_z$ is
invariant under translations of the covariable by every covector
in the span of the differentials of the evaluation $U\mapsto
U n_S(z)$.

For $z\ne w$, both the $z$ and the $w$ evaluation spans contain
the covector
\begin{equation}
 \ell_{z,w}(\dot U)
 =[\dot U n_S(z),U n_S(w)]
 =-[U n_S(z),\dot U n_S(w)].
 \label{eq:locality-shared-covector}
\end{equation}
The equality is the derivative of the pairing-isometry identity.
Thus the symbols of both $d_z$ and $d_w$ are invariant along this
covariable line.  Equation~\eqref{eq:locality-inverse-height}
also gives, on $\mathfrak B$,
\[
 \|L\psi\|\le C_{z,w,L}
                 \big(\|d_z\psi\|+\|d_w\psi\|\big)
\]
for every operator $L$ in the spanning family.  Apply this estimate
to localized oscillatory tests and take the leading-order limit.
Arbitrary compact inner cutoffs and a countable determining set of
linear phases imply the corresponding pointwise symbol estimate. In detail, squaring
the norm bound gives an integral bound by twice the sum of the two
squared symbol norms; arbitrary cutoffs give this bound pointwise for
rational covariables, and polynomial continuity extends it to every
covariable. It bounds the polynomial $\sigma_L(q,\zeta+t\ell_{z,w})$ for all
$t\in\mathbb R$ by a constant independent of $t$.  A vector-valued
polynomial bounded on the real line is constant on that line.
Therefore every principal symbol is invariant along every
$\ell_{z,w}$.

These covectors span the full cotangent space of $\mathcal X$.
Indeed, if a tangent vector $\dot U$ is annihilated by all of them,
then $[\dot U n_S(z),U n_S(w)]=0$ for every distinct pair; by
continuity the same holds for every pair.  Nondegeneracy and the
spanning of the two lift sets give $\dot U=0$.  Hence every
principal symbol is invariant under all covariable translations.
Positive homogeneity makes it zero.  Descend through the finite
orders to obtain order zero.

The compact graph cutoffs and the cores in
Lemma~\ref{lem:locality-two-copy-domain} now give strong commutation
with every first-copy spectral projection, exactly as in the last
paragraph of Lemma~\ref{lem:locality-symbols}.  Second-copy basis
evaluations recover its complete matrix coordinate, so this proves
commutation of the two entire $\mathcal X$ PVMs.
\end{proof}

\begin{proposition}[Isolated-type commutation]
\label{prop:locality-isolated-commutation}
For each isolated type, all angle-height measurements $Q_z$ strongly
commute as $z$ ranges over its source boundary.  Their common law
is the evaluation law of the smooth cone isometries in
Lemma~\ref{lem:locality-cone-maps}; in particular its angular maps
are nonsingular smooth bijections.  The two normal copies sharing
the trunk slot commute with each other at every pair of labels.
All these measurements retain the source and measured-group
covariances, base decomposability, and complementary-group
commutation of the original angle and height measurements.
\end{proposition}

\begin{proof}
Lemmas~\ref{lem:locality-symbols},
\ref{lem:locality-cone-maps}, and
\ref{lem:locality-two-copy-commutation} give the first three
assertions.  Smooth diffeomorphisms of the compact boundaries
preserve the ordinary measure class in both directions.
Every constructed coordinate is obtained by the unique joint
spectral calculus of the original measurements; its finite basis
evaluation recovers that calculus.  Thus their covariance and
commutation identities extend to the joint law, first on a
countable determining algebra and then by normality.  No change
of representation or enlargement of the original commutant is
made in this passage.
\end{proof}
